\documentclass[../mathNotesPreamble]{subfiles}

\providecommand{\relscalefact}{1.4}
\begin{document}
\relscale{\relscalefact}
  \section{8.4: Modular Arithmetic with Applications to Cryptography}
  \begin{defn*}
    The \textbf{Caesar cipher} encrypts messages by shifting letters by a fixed amount (typically 3 places). The encrypted messages are decrypted by shifting in the opposite direction.
  \end{defn*}
  \begin{ex*}
    Use a Ceasar cipher with a shift of $3$ to
    \begin{extasks}[after-item-skip=\stretch{1}](1)
      \task encrypt the message MATH RULES
      \task decrypt the message ODQGHU
    \end{extasks}
    \vspace*{\stretch{1}}
  \end{ex*}
  \pagebreak

  \begin{defn*}
    Given integers $a$ and $n$ with $n>1$, the \textbf{residue of $a$ modulo $n$} is $a \mod n$, the non-negative remainder obtained when $a$ is divided by $n$.
    \begin{itemize}
      \item The numbers $0, 1, \dots, n-1$ are called a \textbf{complete set of residues modulo $n$.}
      \item \textbf{Reducing a number modulo $n$} means to set it equal to its residue modulo $n$.
      \item If a modulus $n>1$ is fixed throughout a discussion and an integer $a$ is given, the words ``modulo $n$'' are often dropped and we simply speak of \textbf{the residue of $a$.}
    \end{itemize}
  \end{defn*}

  \begin{thmBox*}[Theorem 8.4.2: Congruence Modulo $n$ is an Equivalence Relation]
    If $n$ is any integer with $n>1$, congruence modulo $n$ is an equivalence relation on the set of all integers. The distinct equivalence classes of the relation are the sets $\sbrkt{0}, \sbrkt{1}, \sbrkt{2},\dots,\sbrkt{n-1}$, where for each $a=0,1,2,\dots,n-1$,
    \begin{align*}
      \sbrkt{a}&=\set{m\in\bbz \;\middle|\; m\equiv a \mod n}\\
        &=\set{m\in\bbz \;\middle|\; m=a+kn \textnormal{ for some integer } k}
    \end{align*}
  \end{thmBox*}

  \begin{thmBox*}[Theorem 8.4.3: Modular Arithmetic]
    Let $a,b,c,d,$ and $n$ be integers with $n>1$, and suppose
      \[a\equiv c\mod n \qquad \textnormal{ and } \qquad b\equiv d \mod n.\]
    Then
    \begin{enumerate}
      \item $(a+b)\equiv (c+d)\mod n$
      \item $(a-b)\equiv (c-d)\mod n$
      \item $\mathrlap{ab}\phantom{a^m}\equiv \phantom{c^m} \mathllap{cd} \mod n$
      \item $a^m\equiv c^m \mod n$, for every positive integer $m$.
    \end{enumerate}
  \end{thmBox*}
  \pagebreak

  \begin{thmBox*}[Corollary 8.4.4]
    Let $a,b$ and $n$ be integers with $n>1$. Then
      \[ab\equiv\sbrkt{(a\operatorname{\,mod} n)(b\operatorname{\,mod} n)}(\operatorname{mod} n)\]
    or, equivalently,
      \[ab\mod n =\sbrkt{(a\operatorname{\,mod} n)(b\operatorname{\,mod} n)}(\operatorname{mod} n)\]
    In particular, if $m$ is a positive integer, then
      \[a^m \equiv \sbrkt{(a\operatorname{\,mod} n)^m}(\operatorname{mod} n).\]
  \end{thmBox*}

  \begin{ex*}
    Compute the following
    \begin{extasks}[after-item-skip=\stretch{1}](2)
      \task $55+26\mod 4$
      \task $55-26\mod 4$
      \task $55\cdot 26\mod 4$
      \task $55^2\mod 4$
    \end{extasks}
    \vspace*{\stretch{1}}
  \end{ex*}
  \pagebreak

  \noindent
  Recall the following property of exponents:
  \begin{itemize}
    \item $x^{2a}=(x^a)^2$
    \item $x^{a+b}=x^ax^b$
  \end{itemize}
  for all real numbers $x, a,$ and $b$ with $x\geq 0$. By Corollary 8.4.4,
    \[x^7 \mod n = \sbrkt{(x^4 \operatorname{\, mod}n)(x^2 \operatorname{\, mod}n)(x^1 \operatorname{\, mod}n)}\mod n.\]
%  \vspace*{0.5\baselineskip}

  \begin{ex*}
    Compute the following:
    \begin{extasks}[after-item-skip=\stretch{1}](2)
      \task $144^4\mod 713$
      \task $12^{43}\mod 713$
    \end{extasks}
    \vspace*{\stretch{1}}
  \end{ex*}
  \pagebreak

  \begin{defn*}
    An integer $d$ is said to be a \textbf{linear combination of integers} $a$ and $b$ if, and only if, there exists integers $s$ and $t$ such that $as+bt=d$.
  \end{defn*}
%  \pagebreak

  \begin{thmBox*}[Theorem 8.4.5: Writing a Greatest Common Divisor as a Linear Combination]
    For all integers $s$ and $b$, not both zero, if $d=\gcd(a,b)$, then there exist integers $s$ and $t$ such that $as+bt=d$.
  \end{thmBox*}

  \begin{ex*}
    Use the extended Euclidean algorithm to express $\gcd(660,43)$ as a linear combination of $660$ and $43$.
    \vspace*{\stretch{1}}
  \end{ex*}
  \pagebreak

  \begin{defn*}
    Given any integer $a$ and any positive integer $n$, if there exists an integer $s$ such that $as\equiv 1 (\mod n)$, then $s$ is called \textbf{an inverse for $a$ modulo $n$}.

    Integers $a$ and $b$ are \textbf{relatively prime} if, and only if, $\gcd(a,b)=1$.
  \end{defn*}

  \begin{thmBox*}[Corollary 8.4.6]
    If $a$ and $b$ are relatively prime integers, then there exist integers $s$ and $t$ such that $as+bt=1$.
    \vspace*{0.5\baselineskip}

    \noindent
    \textbf{Corollary 8.4.7}

    \noindent
    For all integers $a$ and $n$, if $\gcd(a,n)=1$, then there exists an integer $s$ such that $as\equiv 1 (\mod n)$, and so $s$ is an inverse for $a$ modulo $n$.
  \end{thmBox*}

%  \noindent
%  \textbf{Finding an Inverse Modulo $n$:}
  \begin{ex*}
    \mbox{}
    \begin{extasks}[after-item-skip=\stretch{1}](1)
      \task Find an inverse for $43$ modulo $660$. (Solve $43s\equiv 1 (\mod 660)$.)
      \task Find a positive inverse for $3$ modulo $40$. (Solve $3s\equiv 1 (\mod 40)$.)
    \end{extasks}
    \vspace*{\stretch{1}}
  \end{ex*}
  \pagebreak

%  \noindent
%%  \textbf{RSA Crypotography}
  \begin{defn*}[RSA Cryptography]
    \textbf{RSA Encryption}, named after Rivest, Shamir, and Adleman, uses modular arithmetic to encrypt and decrypt a plaintext message, $M$, into a ciphertext, $C$.

    Given prime integers $p$ and $q$, and a positive integer $e$ that is relatively prime to $(p-1)(q-1)$,
    \begin{itemize}
      \item
         messages are encrypted using the formula:
          \[C = M^e \mod pq.\]
        The number pair $(pq,e)$ is a \textbf{public key}, which is shared.

      \item
        messages are decrypted using the formula:
          \[M = C^d \mod pq\]
        The number pair $(pq,d)$ is a \textbf{private key}. The integer $d$ is a positive inverse to $e$ modulo $(p-1)(q-1)$.
    \end{itemize}
    The security of RSA encryption relies on the difficulty of factoring very large numbers.
  \end{defn*}
  \pagebreak

  \begin{ex*}
    Let $p=5$ and $q=11$. Since $(p-1)(q-1)=40$, we may choose $e=3$, which is relatively prime to $40$.
    \begin{extasks}[after-item-skip=\stretch{1}](1)
      \task Encrypt the plaintext ``HI''
      \task Decrypt the ciphertext ``17 14''
    \end{extasks}
    \vspace*{\stretch{1}}
  \end{ex*}
  \pagebreak

  \begin{algorithm}[H]
    \LinesNotNumbered
    \TabPositions{40mm}
    \SetKwInOut{Input}{Input}
    \SetKwInOut{Output}{Output}
    \Input{$A,B$ [integers with $A>B>0$]}

    $a:=A,\ b:=B,\ s:=1,\ t:=0,\ u:=0,\ v:=1$\\
    \While{$b\neq 0$}{
      $r:=a\operatorname{\,mod} b$,\tab $q:=a\operatorname{\,div} b$\\
      $a:=b$,\tab $b:=r$\\
      $u_{\textnormal{new}}:=s-uq$,\tab $v_{\textnormal{new}}:=t-vq$\\
      $s:=u$,\tab $t:=v$\\
      $u:=u_{\textnormal{new}}$,\tab $v:=v_{\textnormal{new}}$
    }
    $\gcd:=a$\\
    \Output{$\gcd$ [a positive integer], $s, t$ [integers]}
    \caption{Algorithm 8.4.1: Extended Euclidean Algorithm}
  \end{algorithm}
  \pagebreak

\end{document}
